%% kapa: exam=1 round=21 year=2010 period=1 form=A subject=economics subject_ko=경제원론 seq=26 number=66
%% source: Q-Net 감정평가사 기출문제 (cst00309) artlSeq=5024181 file=extracted/1008/5024181/2010년 제21회 감정평가사 1차시험 기출문제/1교시.hwp
%% conversion: kordoc 4.24.51 -> tools/gen1.py (auto); review=proofread-sonnet; flags=none
\begin{question}{66}
주유소에서 매주 휘발유를 甲은 10리터 넣고, 乙은 10,000원 만큼 넣는다. 다음 설명 중 옳은 것을 모두 고른 것은?
\begin{framedtext}
ㄱ. 甲의 휘발유 수요는 가격에 대하여 비탄력적이다. \par
ㄴ. 乙의 경우 휘발유 가격이 10\% 오르면 주유량을 5\% 줄인다. \par
ㄷ. 휘발유 가격이 리터당 1,000원일 때 세금을 부과하면 甲의 조세부담이 乙보다 크다. \par
\end{framedtext}
\begin{choices}
  \item ㄱ
  \item ㄷ
  \item ㄱ, ㄴ
  \item ㄱ, ㄷ
  \item ㄴ, ㄷ
\end{choices}
\end{question}
\answer{4}
